% ╔══════════════════════════════════════════════════════════════════════════╗
% ║  probe.tex --- PLANTILLA MAESTRA DE REPORTE                      ║
% ║  reporte_base.tex  |  Versión 1.0                                       ║
% ║  Motor: Tectonic   |  Preprocesador: Jinja2                             ║
% ╚══════════════════════════════════════════════════════════════════════════╝

\documentclass[11pt, a4paper, final]{article}


% ════════════════════════════════════════════════════════════════════════════
%  0. PAQUETES --- ORDENADOS POR DEPENDENCIA
% ════════════════════════════════════════════════════════════════════════════

% --- Codificación y tipografía ---
\usepackage[T1]{fontenc}
\usepackage[utf8]{inputenc}

% Fuente principal: Helvetica (o clon similar), ultra limpia y corporativa/industrial.
\usepackage{helvet}
\renewcommand{\familydefault}{\sfdefault}

% Fuente secundaria (para \texttt{}): IBM Plex Mono o Inconsolata para los datos crudos.
\usepackage[scale=0.95]{plex-mono} 
\usepackage{microtype}

% --- Idioma (es-noquoting: previene secuestro de < > por babel) ---
\usepackage[spanish, es-tabla, es-noquoting]{babel}

% --- Layout ---
\usepackage[
  a4paper,
  top=2.4cm, bottom=2.6cm,
  left=2.0cm, right=2.0cm,
  headheight=1.6cm,
  headsep=0.6cm,
  footskip=1.0cm
]{geometry}

% --- Color (cargado antes que los demás para que los usen) ---
\usepackage[table, dvipsnames, svgnames]{xcolor}

% --- Gráficos y diagramas ---
\usepackage{tikz}
\usepackage{pgfplots}
\pgfplotsset{compat=1.18}
\usetikzlibrary{
  calc,
  shapes.geometric,
  arrows.meta,
  positioning,
  decorations.pathreplacing,
  backgrounds,
  patterns,
  fit,
  shadows,
  babel
}

% --- Cajas de alerta y contenido ---
\usepackage{tcolorbox}
\tcbuselibrary{skins, breakable, theorems, hooks}

% --- Tablas avanzadas ---
\usepackage{booktabs}
\usepackage{array}
\usepackage{multirow}
\usepackage{tabularx}
\usepackage{colortbl}

% --- MAGIA PARA CENTRAR TABLAS ---
\renewcommand{\tabularxcolumn}[1]{>{\raggedright\arraybackslash}m{#1}}
\renewcommand{\arraystretch}{1.3}

% --- Encabezados y pies ---
\usepackage{fancyhdr}
\usepackage{lastpage}

% --- Matemáticas ---
\usepackage{amsmath}
\usepackage{amssymb}

% --- Utilidades de texto ---
\usepackage{parskip}
\usepackage{enumitem}
\usepackage{graphicx}
\usepackage{calc}
\usepackage{setspace}

% --- Hipervínculos (siempre al final) ---
\usepackage[
  hidelinks,
  pdfauthor={probe.tex},
  pdftitle={Informe de Diagnóstico Forense de Hardware},
  pdfsubject={Hardware Diagnostic Report}
]{hyperref}


% ════════════════════════════════════════════════════════════════════════════
%  1. PALETA DE COLORES CORPORATIVOS
% ════════════════════════════════════════════════════════════════════════════

% --- Escala primaria (azul marino / pizarra) ---
\definecolor{AMPrimary}   {HTML}{000000}
\definecolor{AMSecondary} {HTML}{333333}
\definecolor{AMAccent}    {HTML}{EE0000}
\definecolor{AMSlate}     {HTML}{555555}
\definecolor{AMLight}     {HTML}{F0F0F0}
\definecolor{AMBorder}    {HTML}{DDDDDD}
\definecolor{AMText}      {HTML}{111111}
\definecolor{AMCanvas}    {HTML}{FFFFFF}

% --- Colores de estado (semáforo diagnóstico) ---
\definecolor{StatusOK}     {HTML}{1B4332} % Verde táctico oscuro
\definecolor{StatusOKBg}   {HTML}{F0FFF4} % Blanco con tinte verde muy sutil
\definecolor{StatusOKMid}  {HTML}{2D6A4F} % Verde industrial para gráficos
\definecolor{StatusWarn}   {HTML}{FF5000}
\definecolor{StatusWarnBg} {HTML}{FFF4E6}
\definecolor{StatusWarnMid}{HTML}{FF8C00}
\definecolor{StatusCrit}   {HTML}{EE0000}
\definecolor{StatusCritBg} {HTML}{FDEEED}
\definecolor{StatusCritMid}{HTML}{FF3333}
\definecolor{StatusInfo}   {HTML}{888888}
\definecolor{StatusInfoBg} {HTML}{F5F5F5}
\definecolor{redtex}       {HTML}{FF0000} % Rojo vibrante terminal


% ════════════════════════════════════════════════════════════════════════════
%  2. CONFIGURACIÓN TIPOGRÁFICA Y SECCIONES
% ════════════════════════════════════════════════════════════════════════════

\usepackage[explicit]{titlesec}

% Sección: línea izquierda coloreada + numeración en azul
\titleformat{\section}
  {\large\bfseries\color{AMPrimary}}
  {}
  {0pt}
  {\textcolor{AMPrimary}{\rule[-0.2ex]{4pt}{1.2em}}\hspace{8pt}#1}
\titlespacing*{\section}{0pt}{14pt}{6pt}

% Subsección: línea base punteada
\titleformat{\subsection}
  {\normalsize\bfseries\color{AMSecondary}}
  {\thesubsection.}
  {6pt}
  {#1}
\titlespacing*{\subsection}{0pt}{10pt}{4pt}


% ════════════════════════════════════════════════════════════════════════════
%  3. CAJAS tcolorbox PERSONALIZADAS
% ════════════════════════════════════════════════════════════════════════════
%  Uso en el documento:
%    \begin{amboxok}{Título}   ... \end{amboxok}
%    \begin{amboxwarn}{Título} ... \end{amboxwarn}
%    \begin{amboxcrit}{Título} ... \end{amboxcrit}
%    \begin{amboxinfo}{Título} ... \end{amboxinfo}

% Macro interna para todas las cajas de estado
\newcommand{\amstateboxstyle}[5]{%
  % #1=colback  #2=colframe  #3=symbolo  #4=titulo  #5=contenido
  \begin{tcolorbox}[
    enhanced,
    colback=#1,
    colframe=#2,
    coltitle=#2,
    fonttitle=\bfseries\small,
    title={#3\enspace #4},
    boxrule=0.75pt,
    arc=0pt,
    left=8pt, right=8pt, top=5pt, bottom=5pt
  ]
    #5
  \end{tcolorbox}%
}

% PASS / OK --- Verde -> OPTIMAL
\newtcolorbox{amboxok}[1]{
  enhanced,
  colback=StatusOKBg, colframe=StatusOK,
  coltitle=white, fonttitle=\bfseries\small, % <--- COLTITLE BLANCO
  title={[ OPTIMAL ]\enspace #1},
  boxrule=0.75pt, arc=0pt,
  left=8pt, right=8pt, top=5pt, bottom=5pt
}

% WARNING --- Naranja -> DEGRADED
\newtcolorbox{amboxwarn}[1]{
  enhanced,
  colback=StatusWarnBg, colframe=StatusWarn,
  coltitle=white, fonttitle=\bfseries\small, % <--- COLTITLE BLANCO
  title={[ DEGRADED ]\enspace #1},
  boxrule=0.75pt, arc=0pt,
  left=8pt, right=8pt, top=5pt, bottom=5pt
}

% CRITICAL --- Rojo -> CRITICAL
\newtcolorbox{amboxcrit}[1]{
  enhanced,
  colback=StatusCritBg, colframe=StatusCrit,
  coltitle=white, fonttitle=\bfseries\small, % <--- COLTITLE BLANCO
  title={[ CRITICAL ]\enspace #1},
  boxrule=0.75pt, arc=0pt,
  left=8pt, right=8pt, top=5pt, bottom=5pt
}

% INFO --- Azul neutro -> UNKNOWN
\newtcolorbox{amboxinfo}[1]{
  enhanced,
  colback=StatusInfoBg, colframe=StatusInfo,
  coltitle=white, fonttitle=\bfseries\small, % <--- COLTITLE BLANCO
  title={[ UNKNOWN ]\enspace #1},
  boxrule=0.75pt, arc=0pt,
  left=8pt, right=8pt, top=5pt, bottom=5pt
}

% Caja de info de componente (banda izquierda, sin título flotante)
\newtcolorbox{componentinfo}{
  enhanced,
  colback=AMLight, colframe=AMPrimary,
  boxrule=0pt, leftrule=4pt,
  arc=0pt,
  left=10pt, right=8pt, top=5pt, bottom=5pt
}


% ════════════════════════════════════════════════════════════════════════════
%  4. CONFIGURACIÓN GLOBAL DE pgfplots
% ════════════════════════════════════════════════════════════════════════════

\pgfplotsset{
  % Estilo base aplicable con: \begin{axis}[probe.tex, ...]
  probe.tex/.style={
    % Ejes
    axis line style   = {color=AMSlate, thin},
    tick style        = {color=AMSlate, thin},
    tick label style  = {font=\scriptsize, color=AMText},
    label style       = {font=\footnotesize\bfseries, color=AMSlate},
    title style       = {font=\small\bfseries, color=AMPrimary, yshift=2pt},
    % Grilla
    grid              = both,
    grid style        = {line width=0.25pt, draw=AMBorder},
    major grid style  = {line width=0.4pt,  draw=AMBorder},
    minor grid style  = {line width=0.2pt,  draw=AMBorder!50},
    % Leyenda
    legend style = {
      font         = \scriptsize,
      draw         = AMBorder,
      fill         = AMCanvas,
      rounded corners = 0pt,
      inner sep    = 4pt,
      row sep      = 2pt,
    },
    % Fondo del área de plot
    axis background/.style = {fill=none},
    % Márgenes internos
    enlargelimits = 0.05,
  },
}


% ════════════════════════════════════════════════════════════════════════════
%  5. COMANDOS UTILITARIOS
% ════════════════════════════════════════════════════════════════════════════

% --- Badges de estado inline ---
% Uso: \badgeok  \badgefail  \badgewarn  \badgeinfo
\newcommand{\makebadge}[2]{%
  \tikz[baseline=-0.65ex, inner sep=0pt]\node[ %
    fill=#1, text=white,
    font=\bfseries\scriptsize,
    inner xsep=4pt, inner ysep=2pt,
    rounded corners=0pt
  ]{#2};%
}
\newcommand{\badgeok}  {\makebadge{StatusOK}{OPTIMAL}}
\newcommand{\badgefail}{\makebadge{StatusCrit}{CRITICAL}}
\newcommand{\badgewarn}{\makebadge{StatusWarn}{DEGRADED}}
\newcommand{\badgeinfo}{\makebadge{StatusInfo}{UNKNOWN}}

% --- Valor métrico destacado ---
% Uso: \mv{95$^\circ$C}  ->  valor en azul primario negrita
\newcommand{\mv}[1]{\textcolor{AMPrimary}{\textbf{#1}}}

% --- Separador horizontal de sección ---
\newcommand{\sectionrule}{%
  \vspace{4pt}%
  \noindent\textcolor{AMBorder}{\rule{\linewidth}{0.5pt}}%
  \vspace{4pt}%
}

% --- Fila de tabla coloreada alternada ---
\newcommand{\rowalt}{\rowcolor{AMLight}}

% --- Celda de encabezado de tabla ---
\newcommand{\thcell}[1]{\color{AMPrimary}\textbf{#1}}


% ════════════════════════════════════════════════════════════════════════════
%  6. ENCABEZADOS Y PIE DE PÁGINA (fancyhdr)
% ════════════════════════════════════════════════════════════════════════════

\pagestyle{fancy}
\fancyhf{}
\renewcommand{\headrulewidth}{0pt}
\renewcommand{\footrulewidth}{0pt}

% Encabezado: banda coloreada con regla inferior
\fancyhead[L]{%
  \colorbox{AMPrimary}{%
    \hspace{2pt}%
    \color{white}\small\textbf{probe\textcolor{redtex}{.tex}}%
    \hspace{4pt}%
  }%
  \textcolor{AMPrimary}{\rule[-4pt]{1pt}{14pt}}%
  \hspace{4pt}%
  \color{AMSlate}\small\textit{Diagnóstico Forense de Hardware}%
}
\fancyhead[R]{%
  \small\color{AMSlate}%
  Taller Local Demo \textcolor{AMBorder}{\textbullet}\ 2026-04-12 19:07:44%
}
% Línea bajo encabezado
\renewcommand{\headrule}{%
  \vspace{-4pt}%
  \color{AMPrimary}\rule{\headwidth}{1.2pt}%
}
\renewcommand{\headrulewidth}{1.2pt}

% Pie: ID de reporte + paginación
\fancyfoot[L]{%
  \scriptsize\color{AMSlate}%
  ID:\ \texttt{INV-2026-001} \textbar\ v1.0.0%
}
\fancyfoot[C]{%
  \scriptsize\color{AMSlate}%
  Generado por \textbf{probe\textcolor{redtex}{.tex}} --- Una herramienta de diagnóstico de INVARIANT.%
}
\fancyfoot[R]{%
  \scriptsize\color{AMSlate}%
  Pág.\ \textbf{\thepage}\ de\ \textbf{\pageref{LastPage}}%
}
% Línea gris sobre pie
\renewcommand{\footrule}{%
  \color{AMBorder}\rule{\headwidth}{0.5pt}\vskip 2pt%
}
\renewcommand{\footrulewidth}{0.5pt}


% ════════════════════════════════════════════════════════════════════════════
%  INICIO DEL DOCUMENTO
% ════════════════════════════════════════════════════════════════════════════
\begin{document}
\setlength{\parskip}{6pt}


% ────────────────────────────────────────────────────────────────────────────
%  PORTADA
% ────────────────────────────────────────────────────────────────────────────
\thispagestyle{empty}

\vspace*{1cm}
\noindent
{\color{AMSlate}\scriptsize\bfseries\MakeUppercase{Informe Técnico Confidencial} \hfill ID: \texttt{INV-2026-001}}\\[1ex]
{\color{AMPrimary}\rule{\linewidth}{1.5pt}}\\[4ex]

\noindent
{\color{black}\fontsize{42}{48}\selectfont\bfseries Diagnóstico}\\[1ex]
{\color{black}\fontsize{42}{48}\selectfont\bfseries Forense de Hardware} {\color{StatusCrit}\rule{14pt}{14pt}}\\[3ex]
{\color{AMSlate}\Large Hardware Audit Report}

\vspace{2.5cm}

% --- Bloques de Especificaciones ---
\noindent
\begin{minipage}[t]{0.5\textwidth}
  {\scriptsize\bfseries\color{AMSlate}HARDWARE OBJETIVO}\\[1ex]
  {\Large\bfseries\color{AMPrimary}ASUS\ Vivobook E1504FA}\\[1ex]
  {\small\color{AMSlate}S/N:\ \texttt{R9N0CV12M211383}}
\end{minipage}%
\begin{minipage}[t]{0.5\textwidth}
  {\scriptsize\bfseries\color{AMSlate}ENTORNO DE EJECUCIÓN}\\[1ex]
  {\normalsize\bfseries\color{AMPrimary}probe\textcolor{redtex}{.tex} v1.0.0}\\[1ex]
  {\small\color{AMSlate}Kernel Linux:\ 6.19.10-1-cachyos}
\end{minipage}

\vspace{1.5cm}

\noindent
\begin{minipage}[t]{0.5\textwidth}
  {\scriptsize\bfseries\color{AMSlate}CLIENTE / PROPIETARIO}\\[1ex]
  {\normalsize\bfseries\color{AMPrimary}Taller Local Demo}\\[1ex]
  {\small\color{AMSlate}contacto@cliente.com}\\[0.5ex]
  {\small\color{AMSlate}+54 223 000-0000}
\end{minipage}%
\begin{minipage}[t]{0.5\textwidth}
  {\scriptsize\bfseries\color{AMSlate}DIAGNÓSTICO REALIZADO POR}\\[1ex]
  {\normalsize\bfseries\color{AMPrimary}Invariant Systems}\\[1ex]
  {\small\color{AMSlate}Técnico:\ Admin}\\[0.5ex]
  {\small\color{AMSlate}Fecha:\ 2026-04-12 19:07:44}
\end{minipage}

\vspace*{\fill}

% --- Veredicto Global ---
\noindent
\begin{tikzpicture}
  \node[
    fill=AMCanvas,
    inner sep=18pt,
    text width=\textwidth - 36pt,
    draw=AMPrimary, line width=1.5pt
  ]{%
    \begin{minipage}[c]{0.65\textwidth}
      \small\color{AMText}
      \textbf{\color{AMPrimary}ANÁLISIS ESTRUCTURAL Y TERMODINÁMICO}\\[8pt]
      \textbf{Estado:}\ \badgewarn\\[6pt]
      \scriptsize\color{AMSlate}Degradación física activa detectada en uno o más subsistemas. El hardware opera fuera de su especificación óptima. Se requiere intervención técnica en el próximo ciclo de mantenimiento.
    \end{minipage}%
    \hfill
    \begin{minipage}[c]{0.30\textwidth}
      \raggedleft
      {\scriptsize\bfseries\color{AMSlate}ÍNDICE DE ANOMALÍA ($\Delta A$)}\\[4pt]
      {\fontsize{32}{36}\selectfont\bfseries\color{AMPrimary}20}
    \end{minipage}
  };
\end{tikzpicture}

\vspace{1cm}

% --- Pie de portada ---
\begin{center}
  \large\color{AMSlate}
  Generado por \textbf{probe\textcolor{redtex}{.tex}} --- Una herramienta de diagnóstico de INVARIANT.
\end{center}

\newpage


% ────────────────────────────────────────────────────────────────────────────
%  TABLA DE CONTENIDOS
% ────────────────────────────────────────────────────────────────────────────
{
  \hypersetup{linkcolor=AMSecondary}
  \tableofcontents
}
\newpage


% ════════════════════════════════════════════════════════════════════════════
%  SECCIÓN 1 --- PROCESAMIENTO CENTRAL (CPU)
% ════════════════════════════════════════════════════════════════════════════
\section{Procesamiento Central (CPU)}

\begin{componentinfo}
  \small
  \textbf{Modelo:}\ \mv{AMD Ryzen 3 7320U with Radeon Graphics}\hfill
  \textbf{Núcleos/Hilos:}\ \mv{4/8}\hfill
  \textbf{TDP:}\ \mv{0\ W}\hfill
  \textbf{Socket:}\ FT6
\end{componentinfo}

\vspace{0.5cm}


% ────────────────────────────────────────────────────────────────────────────
%  1.1  Estabilidad Térmica Relativa
% ────────────────────────────────────────────────────────────────────────────
\subsection{Estabilidad Térmica Relativa ($\Delta T_{\mathrm{recovery}}$)}

\noindent % <--- BLOQUEA CUALQUIER SANGRÍA AUTOMÁTICA
\begin{minipage}[t]{0.58\textwidth} % <--- ACHICAMOS UN 2% PARA DAR AIRE AL MEDIO
  \vspace{0pt}
  \begin{tikzpicture}[baseline=(current bounding box.north)]
  \begin{axis}[
    probe.tex,
    width  = \linewidth,
    height = 6.2cm,
    xlabel = {Tiempo (s)},
    ylabel = {Temperatura ($^\circ$C)},
    title  = {Curva de Recuperación Térmica --- CPU},
    xmin=0, xmax=120,
    ymin=30, ymax=105,
    xtick distance=20,
    ytick distance=10,
    extra y ticks      = {68.6},
    extra y tick labels = {\empty}, % <--- SILENCIAMOS EL TEXTO EXTERNO
    extra y tick style  = {
      grid       = major,
      grid style = {dashed, thick, color=StatusWarnMid}
    },
    clip = false,
  ]
    % -- Curva de temperatura medida --
    \addplot[
      color=AMPrimary, very thick, smooth, mark=none
    ] coordinates { (0,64.8) (2,92.0) (4,95.5) (6,96.1) (8,96.2) (11,96.2) (13,96.1) (15,96.1) (17,96.1) (19,96.1) (21,96.2) (23,96.2) (25,96.2) (27,96.2) (30,96.2) (32,90.1) (34,86.2) (36,85.0) (38,87.4) (40,87.4) (42,86.1) (44,82.4) };
    \addlegendentry{$T_{\text{cpu}}$ (medida)}

    % -- Línea de T_idle de referencia --
    \addplot[
      color=StatusOKMid, thick, dashed, mark=none
    ] coordinates { (0,63.6) (120,63.6) };
    \addlegendentry{$T_{\text{idle}}$ = 63.6$^\circ$C}

    % -- Banda de peligro TjMax (zona roja translúcida) --
    \addplot[
      fill=StatusCrit, fill opacity=0.10, draw=none
    ] coordinates { (0,100) (120,100) (120,110) (0,110) }
    \closedcycle;

    % -- Nodo TjMax --
    \node[
      anchor=west, font=\tiny\bfseries, color=StatusCrit
    ] at (axis cs:2, 100+2) {TjMax (100$^\circ$C)};

    % -- NUEVO NODO: T_idle + 5 (Flotando adentro del gráfico) --
    \node[
      anchor=south west, font=\tiny\bfseries, color=StatusWarn
    ] at (axis cs:2, 68.6) {$T_{\text{idle}}+5^\circ$};

  \end{axis}
  \end{tikzpicture}
\end{minipage}%
\hfill
\begin{minipage}[t]{0.39\textwidth} % <--- DAMOS UN 2% MÁS DE ANCHO A LA TABLA
  \vspace{0pt}
  \small
  \begin{tabularx}{\linewidth}{@{}X r@{}}
    \toprule
    \thcell{Parámetro} & \thcell{Valor} \\
    \midrule
    $T_{\max}$ bajo carga        & \mv{96.2$^\circ$C} \\
    \rowalt
    $T_{\text{idle}}$ medida     & \mv{63.6$^\circ$C} \\
    $\Delta T_{\text{recovery}}$ & \mv{32.6$^\circ$C} \\
    \rowalt
    Tiempo de recuperación       & \mv{0.0\ s} \\
    Límite TjMax                 & \mv{100$^\circ$C} \\
    \rowalt
    Pasta térmica (estimada)     & Degradada — sustituir pronto \\
    \bottomrule
  \end{tabularx}

  \vspace{0.5cm}

    \begin{amboxcrit}{Fallo Térmico Detectado}
    \scriptsize $\Delta T$ supera umbral seguro.\\
    Intervención inmediata requerida.
  \end{amboxcrit}
  \end{minipage}


% ────────────────────────────────────────────────────────────────────────────
%  1.2  Frecuencia de Throttling y Estabilidad de P-States
% ────────────────────────────────────────────────────────────────────────────
\subsection{Frecuencia de Throttling y Estabilidad de P-States}

\begin{tabularx}{\linewidth}{@{} X r r r @{}}
  \toprule
  \thcell{Modo de Operación} &
  \thcell{Reloj Base (MHz)} &
  \thcell{Reloj Sostenido (MHz)} &
  \thcell{Desviación} \\
  \midrule
  P0 --- Máximo rendimiento  & \mv{4151} & 3984 & 4.0\% \\
  \rowalt
  P1 --- Balanceado          & \mv{2698} & 2617 & 3.0\% \\
  P2 --- Ahorro de energía   & \mv{1100} & 1089 & 1.0\% \\
  \rowalt
  Throttling térmico (TJ90)  & ---                    & 3113 & --- \\
  \bottomrule
\end{tabularx}

\vspace{0.3cm}

\begin{minipage}[t]{0.48\linewidth}
  \vspace{0pt} % <--- ANCLA INVISIBLE
  \small
  \textbf{Eventos de throttling:}\ \mv{0}\\
  \textbf{Duración acumulada:}\ \mv{0\ ms}\\
  \textbf{Causa principal:}\ Ninguna
\end{minipage}
\hfill
\begin{minipage}[t]{0.48\linewidth}
  \vspace{0pt} % <--- ANCLA INVISIBLE
    \begin{amboxok}{Sin Throttling}
    \scriptsize Frecuencia plena sostenida durante toda la prueba de carga.
  \end{amboxok}
  \end{minipage}

\vspace{0.6cm}


% ────────────────────────────────────────────────────────────────────────────
%  1.3  Machine Check Exceptions (MCE)
% ────────────────────────────────────────────────────────────────────────────
\subsection{Machine Check Exceptions (MCE)}

\begin{amboxok}{Sin Errores MCE Detectados}
  \small
  No se registraron Machine Check Exceptions en el log del kernel.\\[3pt]
  \textbf{Bancos MCE escaneados:}\ 17\quad
  \textbf{Última verificación:}\ 2026-04-12 19:03
\end{amboxok}

\newpage


% ════════════════════════════════════════════════════════════════════════════
%  SECCIÓN 2 --- PROCESAMIENTO GRÁFICO (GPU)
% ════════════════════════════════════════════════════════════════════════════
\section{Procesamiento Gráfico (GPU)}

\begin{componentinfo}
  \small
  \textbf{GPU:}\ \mv{N/A}\hfill
  \textbf{VRAM:}\ \mv{1\ GB}\ (N/A)\hfill
  \textbf{Driver:}\ N/A\hfill
  \textbf{Bus:}\ PCIe\ 4\ x16
\end{componentinfo}

\vspace{0.5cm}


% ────────────────────────────────────────────────────────────────────────────
%  2.1  Delta de Temperatura del Hotspot
% ────────────────────────────────────────────────────────────────────────────
\subsection{Delta de Temperatura del Hotspot ($\Delta T_{\mathrm{hotspot}}$)}

\begin{center}
\begin{tikzpicture}
\begin{axis}[
  probe.tex,
  xbar,
  width  = 0.88\linewidth,
  height = 4.8cm,
  xlabel = {Temperatura ($^\circ$C)},
  symbolic y coords = {Hotspot,Edge},
  ytick  = data,
  yticklabel style = {font=\small\bfseries, color=AMSlate},
  xmin=0, xmax=125,
  xtick distance=20,
  title = {Temperatura GPU: $T_{\text{edge}}$ vs $T_{\text{hotspot}}$},
  nodes near coords,
  nodes near coords style = {
    font=\scriptsize\bfseries, color=AMText,
    xshift=4pt
  },
  nodes near coords align = {horizontal},
  extra x ticks      = {110},
  extra x tick labels = {Límite \mv{110$^\circ$C}},
  extra x tick style  = {
    grid       = major,
    grid style = {dashed, very thick, color=StatusCrit, opacity=0.7},
    tick label style = {
      color=StatusCrit, font=\scriptsize\bfseries,
      rotate=90, anchor=south east, yshift=4pt
    }
  },
  legend pos = south east,
]
  % -- Barra T_edge --
  \addplot[fill=AMSlate, draw=AMSecondary, fill opacity=0.85]
    coordinates { (76.0,Edge) };
  \addlegendentry{$T_{\text{edge}}$ = 76.0$^\circ$C}

  % -- Barra T_hotspot --
  \addplot[fill=AMSecondary, draw=AMPrimary, fill opacity=0.85]
    coordinates { (0.0,Hotspot) };
  \addlegendentry{$T_{\text{hotspot}}$ = 0.0$^\circ$C}

\end{axis}
\end{tikzpicture}
\end{center}

\noindent\small
$\Delta T_{\text{hotspot}} = T_{\text{hotspot}} - T_{\text{edge}} =\ $\mv{-76.0$^\circ$C}
\hspace{1cm}
\badgeok\ Delta dentro del rango normal ($<$20$^\circ$C)

\vspace{0.5cm}


% ────────────────────────────────────────────────────────────────────────────
%  2.2  Integridad de VRAM
% ────────────────────────────────────────────────────────────────────────────
\subsection{Integridad de VRAM}

\begin{tabularx}{\linewidth}{@{} X c c c c @{}}
  \toprule
  \thcell{Prueba de Memoria} & \thcell{Tamaño} &
  \thcell{Patrón} & \thcell{Errores} & \thcell{Estado} \\
  \midrule
  Escritura/Lectura secuencial & \mv{1\ GB} & \texttt{0xDEADBEEF} &
    \mv{0} &
    \badgeok \\
  \rowalt
  Patrón de bits aleatorio     & \mv{1\ GB} & Random &
    \mv{0} &
    \badgeok \\
  Stress combinado (30 s)      & \mv{1\ GB} & Mixed &
    \mv{0} &
    \badgeok \\
  \rowalt
  Conteo ECC corregibles       & --- & --- & \mv{0} &
    \badgeok \\
  \bottomrule
\end{tabularx}

\vspace{0.5cm}


% ────────────────────────────────────────────────────────────────────────────
%  2.3  Estabilidad del Enlace PCIe y Errores AER
% ────────────────────────────────────────────────────────────────────────────
\subsection{Estabilidad del Enlace PCIe y Errores AER}

\begin{minipage}[t]{0.56\linewidth}
  \begin{tabularx}{\linewidth}{@{} X c c @{}}
    \toprule
    \thcell{Parámetro} & \thcell{Soportado} & \thcell{Activo} \\
    \midrule
    Generación PCIe        & Gen4   & \mv{Gen4} \\
    \rowalt
    Ancho de banda (lanes) & x16   & \mv{x16} \\
    Ancho de banda (GB/s)  & 32 GB/s       & \mv{32 GB/s} \\
    \rowalt
    Errores AER corregibles & ---                         & \mv{0} \\
    Errores AER fatales     & ---                         & \mv{0} \\
    \bottomrule
  \end{tabularx}
\end{minipage}
\hfill
\begin{minipage}[t]{0.41\linewidth}
    \begin{amboxok}{PCIe Estable}
    \scriptsize Enlace en capacidad máxima.\\
    Sin errores AER fatales.
  \end{amboxok}
  \end{minipage}

\newpage


% ════════════════════════════════════════════════════════════════════════════
%  SECCIÓN 3 --- ALMACENAMIENTO NO VOLÁTIL (SSD/NVMe)
% ════════════════════════════════════════════════════════════════════════════
\section{Almacenamiento No Volátil (SSD/NVMe)}

\begin{componentinfo}
  \small
  \textbf{Dispositivo:}\ \texttt{/dev/nvme0n1}\hfill
  \textbf{Modelo:}\ \mv{SAMSUNG MZVL4256HBJD-00BTW}\hfill
  \textbf{Capacidad:}\ \mv{256\ GB}\hfill
  \textbf{FW:}\ HXC70W1Q\\
  \textbf{TBW restante:}\ \mv{107.0\ TB}\ de\ 138.9\ TB\hfill
  \textbf{Horas de uso:}\ \mv{7568\ h}
\end{componentinfo}

\vspace{0.5cm}


% ────────────────────────────────────────────────────────────────────────────
%  3.1  Distribución de Latencia de Cola I/O (Histograma Logarítmico)
% ────────────────────────────────────────────────────────────────────────────
%
%  DISEÑO ANTI-CRASH:
%    - ymin=1 : evita log(0) en el eje Y si algún bucket tiene cero ops.
%    - symbolic x coords usa identificadores limpios (b0..b9) para no
%      inyectar sintaxis matematica dentro del argumento de pgfplots.
%    - xticklabels aplica el formato bonito por separado.
%    - Cada coordenada usa guardia Jinja2: si lat_bN == 0 se inyecta 1.
%
\subsection{Distribución de Latencia de Cola I/O (Histograma Logarítmico)}

\begin{center}
\begin{tikzpicture}
\begin{axis}[
  probe.tex,
  width      = 0.94\linewidth,
  height     = 6.5cm,
  ybar,
  bar width  = 14pt,
  xlabel     = {Latencia ($\mu$s)},
  ylabel     = {Operaciones (escala log)},
  ymode      = log,
  ymin       = 1,
  log origin = infty,
  title      = {Histograma de Latencia I/O --- Distribución por Bucket},
  xtick      = data,
  xticklabel style = {font=\scriptsize, rotate=40, anchor=east},
  symbolic x coords = {b0,b1,b2,b3,b4,b5,b6,b7,b8,b9},
  xticklabels = {
    ${<}1$,$1{-}2$,$2{-}4$,$4{-}8$,$8{-}16$,
    $16{-}32$,$32{-}64$,$64{-}128$,$128{-}256$,${>}256$
  },
  nodes near coords,
  nodes near coords style = {
    font=\tiny\bfseries, color=AMText,
    rotate=80, anchor=west
  },
  enlarge x limits = 0.07,
  legend pos = north east,
]
  \addplot[
    fill=AMSlate, draw=AMPrimary, fill opacity=0.82
  ] coordinates {
    (b0,33)
    (b1,40)
    (b2,10)
    (b3,8)
    (b4,60051)
    (b5,131918)
    (b6,17806)
    (b7,34935)
    (b8,6906)
    (b9,2510)
  };
  \addlegendentry{Frecuencia de operaciones}

  % -- Línea de umbral P99 (bucket b7 = 64-128 µs) --
  \addplot[
    color=StatusWarn, dashed, ultra thick, mark=none, forget plot
  ] coordinates {
    (b7, 10) (b7, 1000000)
  };
  \node[
    color=StatusWarn, font=\tiny\bfseries, rotate=90, anchor=south
  ] at (axis cs:b7, 8000) {Umbral P99};

\end{axis}
\end{tikzpicture}
\end{center}

\noindent\small
\textbf{P50:}\ \mv{0.0\ $\mu$s}\quad
\textbf{P95:}\ \mv{0.0\ $\mu$s}\quad
\textbf{P99:}\ \mv{0.0\ $\mu$s}\quad
\textbf{P99.9:}\ \mv{0.0\ $\mu$s}

\vspace{0.5cm}


% ────────────────────────────────────────────────────────────────────────────
%  3.2  WAF y Métricas de Desgaste NVMe (SMART Extended)
% ────────────────────────────────────────────────────────────────────────────
\subsection{WAF y Métricas de Desgaste NVMe (SMART Extended)}

\begin{tabularx}{\linewidth}{@{} X r r c @{}}
  \toprule
  \thcell{Métrica NVMe} & \thcell{Valor Actual} &
  \thcell{Umbral} & \thcell{Estado} \\
  \midrule
  Write Amplification Factor (WAF) & \mv{1.085} & $\leq 3.0$ &
    \badgeok \\
  \rowalt
  LBA escritos (host)         & \mv{31.95\ TB}  & --- & --- \\
  NAND escritas (total)       & \mv{34.67\ TB} & --- & --- \\
  \rowalt
  Bloques NAND defectuosos    & \mv{0}   & $\leq 50$ &
    \badgeok \\
  Spare blocks restantes      & \mv{100} & $\geq 10$ &
    \badgeok \\
  \rowalt
  Errores ECC corregibles     & \mv{0}   & $< 100$ &
    \badgeok \\
  Temperatura máx. registrada & \mv{35.0$^\circ$C} & $< 70^\circ$C &
    \badgeok \\
  \rowalt
  Porcentaje de vida restante & \mv{77\%} & $\geq 20\%$ &
    \badgeok \\
  \bottomrule
\end{tabularx}

\newpage


% ════════════════════════════════════════════════════════════════════════════
%  SECCIÓN 4 --- MEMORIA PRINCIPAL (RAM)
% ════════════════════════════════════════════════════════════════════════════
\section{Memoria Principal (RAM)}

\begin{componentinfo}
  \small
  \textbf{Total:}\ \mv{8\ GB}\hfill
  \textbf{Tipo:}\ \mv{LPDDR5-5500}\hfill
  \textbf{Slots:}\ \mv{2/2}\hfill
  \textbf{Dual Channel:}\ \badgeok\ Activo\end{componentinfo}

\vspace{0.5cm}


% ────────────────────────────────────────────────────────────────────────────
%  4.1  Tasa de Corrección EDAC / Bit Flipping
% ────────────────────────────────────────────────────────────────────────────
\subsection{Tasa de Corrección EDAC / Bit Flipping}

\begin{minipage}[t]{0.54\linewidth}
  \vspace{0pt} % <--- ANCLA INVISIBLE
  \begin{tabularx}{\linewidth}{@{} X c c c @{}}
    \toprule
    \thcell{Módulo DIMM} & \thcell{Tamaño} &
    \thcell{Errores} & \thcell{Estado} \\
    \midrule
        DIMM 0 & 4 GB & 0 & \badgeok \\
    \rowalt
    DIMM 0 & 4 GB & 0 & \badgeok \\
    \bottomrule
    \multicolumn{2}{@{}l}{\scriptsize\textit{Total errores detectados:}} &
    \multicolumn{2}{c@{}}{\mv{0}} \\
  \end{tabularx}

  \vspace{0.4cm}
  \small
  \textbf{Velocidad efectiva:}\ \mv{5500\ MHz}\\
  \textbf{Latencia CAS medida:}\ \mv{0.0\ ns}\\
  \textbf{Método de prueba:}\ \texttt{edac-utils + memtester}
\end{minipage}%
\hfill
\begin{minipage}[t]{0.43\linewidth}
  \vspace{0pt} % <--- ANCLA INVISIBLE
  % -- Gráfico de pastel (donut) --- Integridad de RAM --
  \begin{center}
  \begin{tikzpicture}
    % Ángulo de la porción íntegra (inyectado por Python)
    % start: 90° (12h) en sentido antihorario
    \pgfmathsetmacro{\intAngle}{360.0}
    \pgfmathsetmacro{\failStart}{90 - \intAngle}

    % -- Sector íntegro (verde) --
    \fill[StatusOKMid, opacity=0.9]
      (0,0) -- (90:1.85cm)
      arc[start angle=90, delta angle=-\intAngle, radius=1.85cm]
      -- cycle;

    % -- Sector defectuoso (rojo, si existe) --
    \fill[StatusCritMid, opacity=0.9]
      (0,0) -- (\failStart:1.85cm)
      arc[start angle=\failStart, delta angle=-(360-\intAngle), radius=1.85cm]
      -- cycle;

    % -- Anillo interior (efecto donut) --
    \fill[AMCanvas] (0,0) circle (1.15cm);

    % -- Texto central --
    \node[align=center] at (0,0) {%
      \textcolor{StatusOK}{\bfseries\small 100.0\%}\\[-1pt]
      \textcolor{AMSlate}{\tiny Íntegra}%
    };

    % -- Bordes --
    \draw[AMBorder, line width=0.4pt] (0,0) circle (1.85cm);
    \draw[AMBorder, line width=0.4pt] (0,0) circle (1.15cm);

    % -- Leyenda --
    \node[right, font=\scriptsize, anchor=west] at (2.05cm,  0.45cm)
      {\textcolor{StatusOKMid}{$\blacksquare$}\ Íntegra\ (100.0\%)};
    \node[right, font=\scriptsize, anchor=west] at (2.05cm, -0.10cm)
      {\textcolor{StatusCritMid}{$\blacksquare$}\ Defectos\ (0.0\%)};

    % -- Título --
    \node[above, font=\small\bfseries, color=AMPrimary]
      at (0, 2.15cm) {Integridad de RAM};
  \end{tikzpicture}
  \end{center}
  \end{minipage}

  \vspace{0.6cm}

    \begin{amboxok}{RAM Sin Errores}
    \scriptsize Todos los módulos DIMM pasaron las pruebas de integridad.
  \end{amboxok}
  
\newpage


% ════════════════════════════════════════════════════════════════════════════
%  SECCIÓN 5 --- PLACA BASE Y SISTEMA DE ENERGÍA
% ════════════════════════════════════════════════════════════════════════════
\section{Placa Base y Sistema de Energía}

\begin{componentinfo}
  \small
  \textbf{Placa Base:}\ \mv{ASUSTeK COMPUTER INC.\ E1504FA}\hfill
  \textbf{Chipset:}\ Advanced Micro Devices, Inc. [AMD] Family 17h-19h PCIe Root Complex\hfill
  \textbf{BIOS/UEFI:}\ E1504FA.314\ (04/11/2025)
\end{componentinfo}

\vspace{0.5cm}


% ────────────────────────────────────────────────────────────────────────────
%  5.1  Caída de Voltaje bajo Carga (V_droop) y VRM
% ────────────────────────────────────────────────────────────────────────────
\subsection{Caída de Voltaje bajo Carga ($V_{\mathrm{droop}}$) y VRM}

\begin{tikzpicture}
\begin{axis}[
  probe.tex,
  width  = 0.92\linewidth,
  height = 5.8cm,
  xlabel = {Tiempo (s)},
  ylabel = {Voltaje (V)},
  title  = {VRM: Voltaje Solicitado (VID) vs Voltaje Entregado (Medido)},
  xmin=0, xmax=60,
  ymin=0.80, ymax=1.55,
  xtick distance=10,
  ytick distance=0.05,
  legend pos=south east,
]
  % -- Zona de tolerancia ±5% (relleno verde translúcido) --
  \addplot[
    fill=StatusOKMid, fill opacity=0.07, draw=none, forget plot
  ] coordinates {
    (0,  1.273)  (60, 1.273)
    (60, 1.407) (0,  1.407)
  } \closedcycle;

  % -- Voltaje solicitado (VID) --
  \addplot[
    color=AMSlate, thick, dashed, mark=none
  ] coordinates { (0,1.34) (2,1.34) (4,1.34) (6,1.3407) (8,1.342) (10,1.3433) (12,1.3447) (14,1.346) (16,1.3473) (18,1.3487) (20,1.35) (22,1.348) (24,1.348) (26,1.348) (28,1.348) (30,1.348) (32,1.348) (34,1.348) (36,1.348) (38,1.348) (40,1.348) (42,1.348) (44,1.348) (46,1.3475) (48,1.3464) (50,1.3453) (52,1.3443) (54,1.3432) (56,1.3421) (58,1.3411) (60,1.34) };
  \addlegendentry{$V_{\text{solicitado}}$ (VID)}

  % -- Voltaje entregado (medido por sensor) --
  \addplot[
    color=AMPrimary, thick, smooth, mark=none
  ] coordinates { (0,1.34) (2,1.34) (4,1.34) (6,1.3383) (8,1.3348) (10,1.3313) (12,1.3279) (14,1.3244) (16,1.3209) (18,1.3175) (20,1.314) (22,1.3145) (24,1.3148) (26,1.315) (28,1.3149) (30,1.3146) (32,1.3141) (34,1.3136) (36,1.3132) (38,1.313) (40,1.313) (42,1.3133) (44,1.3138) (46,1.3157) (48,1.3192) (50,1.3227) (52,1.3261) (54,1.3296) (56,1.3331) (58,1.3365) (60,1.34) };
  \addlegendentry{$V_{\text{entregado}}$ (sensor)}

  % -- Anotación del droop máximo --
  \draw[->, color=AMSlate, thin]
    (axis cs:20.0, 1.314)
    -- (axis cs:20.0+8, 1.314-0.04)
    node[right, font=\tiny, color=AMSlate] {%
      $V_{\text{droop}}$\mv{= 0.026\ V}%
    };

\end{axis}
\end{tikzpicture}

\vspace{0.3cm}
\noindent\small
\textbf{$V_{\text{droop}}$ máx.:}\ \mv{0.026\ V}\quad
\textbf{Temperatura VRM:}\ \mv{0.0$^\circ$C}\quad
\textbf{Fases VRM:}\ \mv{0}\quad
\textbf{Estado:}\ \mv{Dentro de tolerancia}

\vspace{0.3cm}
\begin{amboxok}{VRM Estable}
  \scriptsize Caída de voltaje dentro del margen de tolerancia ($\pm$5\%). Sin riesgo de inestabilidad.
\end{amboxok}

\vspace{0.6cm}


% ────────────────────────────────────────────────────────────────────────────
%  5.2  Estado de Batería (condicional --- solo laptops)
% ────────────────────────────────────────────────────────────────────────────
\subsection{Estado de Batería (si aplica)}


\begin{minipage}[t]{0.54\linewidth}
  \begin{tabularx}{\linewidth}{@{} X r @{}}
    \toprule
    \thcell{Parámetro de Batería} & \thcell{Valor} \\
    \midrule
    Capacidad de diseño         & \mv{45006\ mWh} \\
    \rowalt
    Capacidad actual (plena)    & \mv{42128\ mWh} \\
    Ciclos de carga             & \mv{0} \\
    \rowalt
    Estado de salud (SoH)       & \mv{94\%} \\
    Voltaje nominal             & \mv{11.85\ V} \\
    \rowalt
    Voltaje bajo carga          & \mv{12.678\ V} \\
    $\Delta V$ bajo carga       & \mv{0.0\ V} \\
    \rowalt
    Corriente de descarga       & \mv{0\ mA} \\
    Resistencia interna est.    & \mv{0.0\ m$\Omega$} \\
    \bottomrule
  \end{tabularx}
\end{minipage}%
\hfill
\begin{minipage}[t]{0.43\linewidth}
  % -- Gauge de SoH (barra de rango TikZ) --
  \begin{center}
  \begin{tikzpicture}
    % Constantes de la barra
    \def\bw{4.6}    % ancho total (cm)
    \def\bh{0.6}    % alto (cm)

    % Fondo gris
    \fill[AMBorder!60] (0,0) rectangle (\bw, \bh);

    % Relleno de color según SoH (longitud pre-calculada por Python)
    \fill[StatusOKMid]
      (0,0) rectangle (4.32, \bh);

    % Divisores de zona (0-40%=critico, 40-70%=moderado, 70-100%=bueno)
    \draw[white, line width=1.2pt] (1.84, 0) -- (1.84, \bh);  % 40%
    \draw[white, line width=1.2pt] (3.22, 0) -- (3.22, \bh);  % 70%

    % Indicador de posición actual
    \draw[AMPrimary, very thick]
      (4.32, -0.18) -- (4.32, \bh+0.18);
    % Triángulo indicador
    \fill[AMPrimary]
      (4.32, \bh+0.18) --
      ({4.32-0.12}, \bh+0.38) --
      ({4.32+0.12}, \bh+0.38) -- cycle;

    % Etiquetas de zona
    \node[font=\tiny, color=StatusCrit]
      at (0.92, -0.28) {Crítico};
    \node[font=\tiny, color=StatusWarn]
      at (2.53, -0.28) {Moderado};
    \node[font=\tiny, color=StatusOK]
      at (3.91, -0.28) {Bueno};

    % Borde
    \draw[AMSlate, line width=0.5pt] (0,0) rectangle (\bw, \bh);

    % Valor de SoH
    \node[above, font=\bfseries\small, color=AMPrimary]
      at (\bw/2, \bh+0.55) {Estado de Salud (SoH): \mv{94\%}};

    % Título
    \node[above, font=\scriptsize, color=AMSlate]
      at (\bw/2, \bh+1.0) {Resistencia interna est.: \mv{0.0\ m$\Omega$}};
  \end{tikzpicture}
  \end{center}
  \end{minipage}

  \vspace{0.6cm}

    \begin{amboxok}{Batería en Buen Estado}
    \scriptsize SoH: \mv{94\%}. No requiere reemplazo.
  \end{amboxok}
  

\newpage


% ════════════════════════════════════════════════════════════════════════════
%  SECCIÓN 6 --- INTEGRIDAD DE BUS Y PERIFÉRICOS
% ════════════════════════════════════════════════════════════════════════════
\section{Integridad de Bus y Periféricos}


% ────────────────────────────────────────────────────────────────────────────
%  6.1  Integridad USB
% ────────────────────────────────────────────────────────────────────────────
\subsection{Integridad de Puertos USB}

\begin{tabularx}{\linewidth}{@{} c c X c c c @{}}
  \toprule
  \thcell{Puerto} & \thcell{Versión} & \thcell{Controlador} &
  \thcell{BW (Gbps)} & \thcell{Errores} & \thcell{Estado} \\
  \midrule
      Bus 1 Puerto 002 & USB 2.0 (HS) & xHCI & 0.48 & 0 & \badgeok \\
    \rowalt
    Bus 1 Puerto 004 & USB 2.0 (HS) & xHCI & 0.48 & 0 & \badgeok \\
    Bus 1 Puerto 004 & USB 2.0 (HS) & xHCI & 0.48 & 0 & \badgeok \\
    \rowalt
    Bus 1 Puerto 004 & USB 2.0 (HS) & xHCI & 0.48 & 0 & \badgeok \\
    Bus 3 Puerto 002 & USB 2.0 (HS) & xHCI & 0.48 & 0 & \badgeok \\
    \rowalt
    Bus 3 Puerto 002 & USB 2.0 (HS) & xHCI & 0.48 & 0 & \badgeok \\
    Bus 3 Puerto 002 & USB 2.0 (HS) & xHCI & 0.48 & 0 & \badgeok \\
    \rowalt
    Bus 3 Puerto 003 & USB 1.1 (FS) & xHCI & 0.012 & 0 & \badgeok \\
    Bus 3 Puerto 003 & USB 1.1 (FS) & xHCI & 0.012 & 0 & \badgeok \\
    \rowalt
    Bus 3 Puerto 003 & USB 1.1 (FS) & xHCI & 0.012 & 0 & \badgeok \\
    Bus 5 Puerto 001 & USB 2.0 (HS) & xHCI & 0.48 & 0 & \badgeok \\
    \rowalt
    Bus 5 Puerto 001 & USB 2.0 (HS) & xHCI & 0.48 & 0 & \badgeok \\
    Bus 5 Puerto 001 & USB 2.0 (HS) & xHCI & 0.48 & 0 & \badgeok \\
  \bottomrule
\end{tabularx}

\vspace{0.5cm}

\begin{minipage}[t]{0.48\linewidth}
  \vspace{0pt} % <--- ANCLA INVISIBLE
  \small\textbf{\color{AMPrimary}Resumen de Estado USB:}
  \begin{itemize}[leftmargin=1.2em, itemsep=3pt, topsep=5pt]
    \item \badgeok\ Puertos estables: \mv{13}
    \item \badgewarn\ Puertos inestables: \mv{0}
    \item \badgefail\ Puertos con fallo/corto: \mv{0}
  \end{itemize}
\end{minipage}
\hfill
\begin{minipage}[t]{0.48\linewidth}
  \vspace{0pt} % <--- ANCLA INVISIBLE
    \begin{amboxok}{Todos los Puertos USB Operativos}
    \scriptsize Sin errores detectados en ningún puerto USB.
  \end{amboxok}
  \end{minipage}


% ════════════════════════════════════════════════════════════════════════════
%  TABLA RESUMEN GLOBAL + RECOMENDACIONES
% ════════════════════════════════════════════════════════════════════════════
\vspace{0.8cm}
\sectionrule

\begin{center}
  {\color{AMPrimary}\large\bfseries Resumen Global del Diagnóstico}
\end{center}
\sectionrule

\vspace{0.2cm}

\begin{tabularx}{\linewidth}{@{} X c c l @{}}
  \toprule
  \thcell{Sub-Sistema} & \thcell{$\Delta A$ Parcial} &
  \thcell{Estado Clínico} & \thcell{Anotación} \\
  \midrule
  Procesamiento Central (CPU)   & \textbf{+20}  & \badgewarn  & --- \\
  \rowalt
  Procesamiento Gráfico (GPU)   & \textbf{+0}  & \badgeok  & --- \\
  Almacenamiento NVMe/SSD       & \textbf{+0} & \badgeok & --- \\
  \rowalt
  Memoria Principal (RAM)       & \textbf{+0}  & \badgeok  & --- \\
  Placa Base / VRM Energía      & \textbf{+0} & \badgeok & --- \\
  \rowalt
  Bus / Periféricos (USB)       & \textbf{+0}  & \badgeok  & --- \\
  \rowalt
  Celdas de Energía (BAT)       & \textbf{+0}  & \badgeok  & --- \\
  \midrule
  \thcell{ÍNDICE DE ANOMALÍA GLOBAL} &
  \thcell{20} &
  \thcell{\badgewarn} &
  \thcell{---} \\
  \bottomrule
\end{tabularx}

\vspace{0.6cm}

\begin{tcolorbox}[
  enhanced,
  colback   = AMLight,
  colframe  = AMPrimary,
  boxrule   = 0.8pt,
  arc       = 0pt,
  title     = {\bfseries Directivas de Intervención},
  coltitle  = white,
  attach boxed title to top left={yshift=-2mm, xshift=8mm},
  boxed title style={colback=AMPrimary, arc=0pt, boxrule=0pt},
  left=10pt, right=10pt, top=12pt, bottom=8pt
]
  \small
  \begin{enumerate}[leftmargin=*, itemsep=4pt, topsep=2pt]
    \item \textbf{Sub-sistema CPU / Disipación:} $\Delta T_{\text{recovery}}$ = 32.6°C (umbral: 15°C), $t_{\text{recovery}}$ = 0.0 s (umbral: 45 s). Capacidad de disipación comprometida. Reemplazar pasta térmica. Verificar flujo de aire y disipador.
  \end{enumerate}
\end{tcolorbox}


% ════════════════════════════════════════════════════════════════════════════
%  PIE DE DOCUMENTO --- FIRMA Y SELLO TÉCNICO
% ════════════════════════════════════════════════════════════════════════════
\vspace{0.8cm}
\sectionrule

\begin{minipage}[t]{0.48\linewidth}
  \small\color{AMSlate}
  \textbf{Emitido por:}\ Invariant Systems\\
  \textbf{Técnico certificado:}\ Admin\\
  \textbf{Fecha de diagnóstico:}\ 2026-04-12 19:07:44\\
  \textbf{Duración del análisis:}\ 302.0 s
\end{minipage}
\hfill
\begin{minipage}[t]{0.48\linewidth}
  \raggedright\small\color{AMSlate}
  \textbf{Plataforma:}\ probe\textcolor{redtex}{.tex} v1.0.0\\
  \textbf{Kernel Linux:}\ 6.19.10-1-cachyos\\
  \textbf{ID de Reporte:}\ \texttt{INV-2026-001}\\
  \textbf{Hash de integridad (SHA-256):}\ \texttt{\tiny }
\end{minipage}

\vspace{0.5cm}

\begin{center}
  \scriptsize\color{AMSlate}
  Este informe es emitido con fines diagnósticos y de referencia técnica exclusivamente.\\
  Los datos fueron recolectados de forma automatizada por probe\textcolor{redtex}{.tex}
  en un entorno Live Linux offline controlado.\\[3pt]
  \textcolor{AMBorder}{\rule{0.45\linewidth}{0.3pt}}
\end{center}

\end{document}